% ============================================================
%  \subsection{Evaluation metrics}
%  Drop this into main.tex under Section: Identity Evaluation
% ============================================================

\subsection{Evaluation metrics}
\label{subsec:metrics}

We formulate identity verification as a binary classification task
and evaluate it through a multi-tier framework that progressively
captures both recognition accuracy and end-to-end system reliability.

%------------------------------------------------------------
\subsubsection{Binary verification formulation}
\label{subsubsec:binary}

Given a cosine-similarity score $s = \operatorname{sim}(\mathbf{e}_a, \mathbf{e}_b)$
and a decision threshold $\tau$, we classify each pair as:
\begin{equation}
    \hat{y} =
    \begin{cases}
        1 & \text{if } s \geq \tau \quad \text{(same identity)}, \\
        0 & \text{otherwise} \quad \text{(different identity)}.
    \end{cases}
    \label{eq:decision}
\end{equation}

The standard confusion-matrix elements are defined accordingly:
\begin{itemize}
    \item \textbf{True Positive (TP)}: $\hat{y} = 1$ and $y = 1$ (correctly matched same identity).
    \item \textbf{True Negative (TN)}: $\hat{y} = 0$ and $y = 0$ (correctly rejected different identity).
    \item \textbf{False Positive (FP)}: $\hat{y} = 1$ and $y = 0$ (falsely matched different identity).
    \item \textbf{False Negative (FN)}: $\hat{y} = 0$ and $y = 1$ (falsely rejected same identity).
\end{itemize}

From these elements, we compute four primary metrics:
\begin{align}
    \text{Accuracy}  &= \frac{|\text{TP}| + |\text{TN}|}{|\text{TP}| + |\text{TN}| + |\text{FP}| + |\text{FN}|}, \label{eq:acc} \\[4pt]
    \text{Precision} &= \frac{|\text{TP}|}{|\text{TP}| + |\text{FP}|}, \label{eq:prec} \\[4pt]
    \text{Recall}    &= \frac{|\text{TP}|}{|\text{TP}| + |\text{FN}|}, \label{eq:rec} \\[4pt]
    F_1              &= \frac{2 \cdot \text{Precision} \cdot \text{Recall}}
                           {\text{Precision} + \text{Recall}}
                      = \frac{2\,|\text{TP}|}
                             {2\,|\text{TP}| + |\text{FP}| + |\text{FN}|}. \label{eq:f1}
\end{align}

%------------------------------------------------------------
\subsubsection{Balanced pair construction}
\label{subsubsec:pairs}

To eliminate class-imbalance bias, we construct a \emph{1:1 balanced}
evaluation set for each experimental condition $(v, \sigma)$, where
$v \in \{\text{frontal}, \text{side}\}$ is the view type and $\sigma$
is the blur level.

For $N$ subjects in a given condition, we define:
\begin{itemize}
    \item \textbf{Positive pairs} ($\mathcal{P}^{+}$): For each subject $i$,
          pair the restored/degraded embedding $\mathbf{e}_{\text{test}}^{(i)}$
          with its own reference $\mathbf{e}_{\text{ref}}^{(i)}$:
    \begin{equation}
        \mathcal{P}^{+} = \bigl\{
            (\mathbf{e}_{\text{test}}^{(i)},\;
             \mathbf{e}_{\text{ref}}^{(i)},\;
             y{=}1)
            \;\big|\; i = 1, \dots, N
        \bigr\}.
        \label{eq:pos_pairs}
    \end{equation}

    \item \textbf{Negative pairs} ($\mathcal{P}^{-}$): For each subject $i$,
          pair the restored/degraded embedding with a randomly selected
          reference from a \emph{different} subject $j \neq i$:
    \begin{equation}
        \mathcal{P}^{-} = \bigl\{
            (\mathbf{e}_{\text{test}}^{(i)},\;
             \mathbf{e}_{\text{ref}}^{(j)},\;
             y{=}0)
            \;\big|\; j \sim \mathcal{U}\bigl(\{1,\dots,N\} \setminus \{i\}\bigr)
        \bigr\}.
        \label{eq:neg_pairs}
    \end{equation}
\end{itemize}

\noindent
This yields $|\mathcal{P}^{+}| = |\mathcal{P}^{-}| = N$ pairs,
forming a balanced dataset of $2N$ total pairs per condition.
A fixed random seed ($\text{seed}=42$) ensures reproducibility of
the negative-pair assignment across experiments.

%------------------------------------------------------------
\subsubsection{Dual threshold strategy}
\label{subsubsec:threshold}

We report results under two complementary thresholding protocols:

\paragraph{Fixed threshold ($\tau_{\text{fix}} = 0.15$).}
This pre-defined operating point is chosen for degraded/restored
face scenarios, where similarity scores are substantially lower
than in clean-image face verification (where thresholds typically
exceed 0.40). Using a fixed threshold enables direct cross-condition
and cross-method comparability without overfitting to a specific
similarity distribution.

\paragraph{Optimal threshold ($\tau^{*}$).}
We perform an exhaustive grid search over
$\tau \in \{-0.50,\; -0.49,\; \dots,\; 0.99\}$ and select the threshold
that maximizes the $F_1$-score:
\begin{equation}
    \tau^{*} = \arg\max_{\tau} \; F_1(\tau).
    \label{eq:optimal_threshold}
\end{equation}

The optimal threshold reveals the \emph{best achievable discrimination}
of a given method and indicates how well-separated the positive and
negative similarity distributions are. A high optimal $F_1$ with
$\tau^{*} \gg 0$ reflects strong identity preservation, whereas
$\tau^{*} \leq 0$ indicates that positive and negative distributions
overlap severely, signaling near-complete identity loss.

%------------------------------------------------------------
\subsubsection{Three-tier evaluation hierarchy}
\label{subsubsec:tiers}

We define three evaluation tiers to progressively disentangle
face detection reliability from recognition quality.

\paragraph{Tier~1: Face Detection Rate (FDR).}
The proportion of input images (or video frames) for which the
face recognition system successfully detects and embeds a face:
\begin{equation}
    \text{FDR} = \frac{N_{\text{detected}}}{N_{\text{total}}} \times 100\%.
    \label{eq:fdr}
\end{equation}

Under severe degradation, face detectors may fail entirely on
blurred images while still detecting faces in video-generated
frames. This tier captures a critical operational dimension: a
restoration method that produces undetectable faces has zero
downstream utility, regardless of its pixel-level quality.

\paragraph{Tier~2: Intersection metrics.}
Computed \emph{only} on the subset of pairs where both
the test and reference faces were successfully detected:
\begin{equation}
    \mathcal{S}_{\text{valid}} = \bigl\{
        (a, b, y) \in \mathcal{P}
        \;\big|\;
        \operatorname{sim}(a, b) \neq \texttt{NaN}
    \bigr\}.
    \label{eq:valid}
\end{equation}

Metrics at this tier answer: \emph{``When a face is detected, how
accurately is its identity preserved?''} By excluding detection
failures, Tier~2 isolates the intrinsic recognition quality of each
method from its detection robustness.

\paragraph{Tier~3: System-level metrics.}
All pairs are included; failed detections ($\texttt{NaN}$ cosine
similarities) are penalized by assigning a score of $-1.0$:
\begin{equation}
    \tilde{s}_{(a,b)} =
    \begin{cases}
        \operatorname{sim}(\mathbf{e}_a, \mathbf{e}_b) & \text{if both faces detected}, \\
        -1.0 & \text{otherwise}.
    \end{cases}
    \label{eq:penalty}
\end{equation}

The penalty value $-1.0$ is chosen to lie strictly below all
observed genuine similarities (which are $> -0.5$ in practice),
ensuring that any threshold $\tau > -1.0$ classifies a detection
failure as a rejected match. This tier captures the
\emph{end-to-end operational performance}, combining both detection
reliability and recognition accuracy into a single set of metrics.
It is the most operationally relevant measure for real-world
deployment, where a pipeline failure at any stage results in a
verification failure.

%------------------------------------------------------------
\subsubsection{Dual face recognition systems}
\label{subsubsec:dual_fr}

To ensure that our findings generalize beyond a single recognition
backend, we evaluate each condition with two independent face
recognition systems that share the ArcFace~\cite{ArcFace} backbone
but differ in detector, preprocessing, and pose estimation:

\begin{enumerate}
    \item \textbf{InsightFace (buffalo\_l)}: Uses
          SCRFD~\cite{SCRFD} for face detection, ArcFace for 512-d
          embedding extraction, and the Simple3D module for direct
          3D head-pose estimation ($\psi, \theta, \phi$).
    \item \textbf{DeepFace (ArcFace + RetinaFace)}: Uses
          RetinaFace~\cite{RetinaFace} for detection with affine
          alignment, ArcFace via the DeepFace library for embedding
          extraction, and a landmark-based proxy for frontality
          estimation (see Section~\ref{subsec:frame_selection}).
\end{enumerate}

The two systems exhibit markedly different behaviors under
degradation: InsightFace's SCRFD detector is more conservative
and may fail to detect severely blurred faces, whereas DeepFace's
RetinaFace (with \texttt{enforce\_detection=False}) can produce
embeddings even from heavily degraded inputs, potentially at the
cost of embedding quality. Reporting results from both systems
strengthens our conclusions through:
\begin{enumerate*}[label=(\roman*)]
    \item cross-validation of the I2V advantage,
    \item detector-sensitivity analysis, and
    \item face-detection-rate comparison.
\end{enumerate*}

%------------------------------------------------------------
\subsubsection{Experimental conditions matrix}
\label{subsubsec:conditions}

The full factorial experimental design is summarized in
Table~\ref{tab:conditions}.

\begin{table}[h]
\centering
\caption{Experimental conditions matrix.}
\label{tab:conditions}
\begin{tabular}{l l c}
\toprule
\textbf{Dimension} & \textbf{Values} & \textbf{Count} \\
\midrule
View type ($v$) & Frontal, Side & 2 \\
Blur level ($\sigma$) & 3,\,5,\,8,\,10,\,12,\,15 & 6 \\
Restoration method & I2V, SR, Degraded & 3 \\
Recognition system & InsightFace, DeepFace & 2 \\
\midrule
\multicolumn{2}{l}{\textbf{Total unique conditions}} & \textbf{72} \\
\bottomrule
\end{tabular}
\end{table}

\noindent
Each condition is evaluated on $2N = 500$ pairs (250~positive +
250~negative), yielding a total of $\mathbf{36{,}000}$ pair-wise
identity verification comparisons across the entire study.
